% Final
\chapter{Návrh}

% Final
Výstupem předchozí kapitoly byl seznam požadavků, které definují jednotlivé funkce a vlastnosti aplikace. Cílem této kapitoly bude na základě těchto požadavků provést návrh a podrobnější specifikaci jednotlivých funkcí a vytvořit tak podrobné podklady pro samotný vývoj aplikace.

% Final
V rámci této kapitoly se nejprve zaměřím na podrobnou specifikaci jednotlivých funkcí aplikace pomocí případů užití a metodiky Behavior-Driven Development \cite{BDD} s využitím jazyka Gherkin \cite{CucumberDocs}. Dále se budu věnovat návrhu samotného vzhledu uživatelského rozhraní včetně výběru palety barev, písma, obrázků, animací a efektů. Následně s využitím tohoto návrhu vytvořím \textit{high fidelity} prototyp uživatelského rozhraní aplikace.

% Final
Díky detailnímu návrhu uživatelského rozhraní a podrobné specifikaci jednotlivých funkcí tak bude aplikace připravena na samotný vývoj, kterému se budu věnovat v následující kapitole.
